% main.tex --- The Collatz Conjecture: Survey, a Natural-Density Obstruction,
%               and Computational Notes.
% Author: Alex Ye. AI assistance (Claude, Anthropic) disclosed via \thanks, not authorship.
%
% VERIFICATION / HONESTY POSTURE (see also the \thanks and Section "Verification status"):
%  - This is an EXPOSITORY survey + one NEGATIVE result (a structural obstruction).
%    It contains NO proof of the Collatz conjecture and NO claim that natural density
%    fails -- only that a specific reduction (the scalar-Esscher mixing route) is obstructed.
%  - Primary-source PDFs (arXiv, publisher) returned HTTP 403 in the working environment.
%    Every external constant/attribution is therefore snippet-sourced and flagged
%    "% [CONSTANT UNVERIFIED vs primary]". Novelty claims are flagged
%    "% [PRIOR-ART CHECK PENDING]" and phrased defensively.

\documentclass[11pt]{amsart}

%--- Encoding & fonts ---
\usepackage[utf8]{inputenc}
\usepackage[T1]{fontenc}
\usepackage{lmodern}

%--- Math packages ---
\usepackage{amsmath, amssymb, amsthm, mathtools}
\usepackage{bm}

%--- Layout & utilities ---
\usepackage[margin=1in]{geometry}
\usepackage{hyperref}
\usepackage{xcolor}
\usepackage{microtype}
\usepackage{enumitem}
\usepackage{booktabs}

%--- Bibliography ---
\usepackage[numbers,sort&compress]{natbib}

%--- Local macros ---
% macros.tex --- Notation and shortcuts for the Collatz paper
% Authors: Alex Ye with Claude (Anthropic)

%--- Lagarias-convention operators ---
\DeclareMathOperator{\ord}{ord}
\DeclareMathOperator{\lcm}{lcm}
\DeclareMathOperator{\val}{v}             % 2-adic valuation
\DeclareMathOperator{\sgn}{sgn}

%--- The Collatz map and iterates ---
\newcommand{\T}{T}                          % Collatz / Syracuse map
\newcommand{\Tit}[1]{T^{#1}}                % T^k
\newcommand{\Cmap}{C}                       % accelerated Syracuse: C(n) = (3n+1)/2^{v_2(3n+1)}

%--- Stopping times ---
\newcommand{\stop}{\sigma}                  % stopping time (first n -> below n)
\newcommand{\totstop}{\sigma_\infty}        % total stopping time (first reach 1)
\newcommand{\stoptime}[1]{\stop(#1)}
\newcommand{\totstoptime}[1]{\totstop(#1)}

%--- Parity vector / Q-functions ---
\newcommand{\pv}{Q}                         % parity vector / Q-function
\newcommand{\pvk}[2]{Q_{#1}(#2)}            % first k entries of parity vector for n
\newcommand{\trajectory}[1]{\mathcal{O}(#1)}% forward orbit

%--- Density / measure ---
\newcommand{\logdens}{\overline{d_{\log}}}  % logarithmic density (upper)
\newcommand{\logdenslow}{\underline{d_{\log}}}
\newcommand{\nat}{\mathbb{N}}
\newcommand{\Z}{\mathbb{Z}}
\newcommand{\Zp}{\mathbb{Z}_p}              % p-adic integers
\newcommand{\Q}{\mathbb{Q}}
\newcommand{\R}{\mathbb{R}}

%--- Tao framework ---
\newcommand{\TaoConst}{c}                   % the implicit constant in Tao 2019
\newcommand{\logbound}{f(N)}                % function s.t. orbits go below f(N)

%--- Workflow markers ---
\newcommand{\TODO}[1]{\textcolor{red}{\textbf{[TODO: #1]}}}
\newcommand{\NOTE}[1]{\textcolor{blue}{\textbf{[NOTE: #1]}}}

%--- Misc shorthands ---
\newcommand{\eps}{\varepsilon}
\newcommand{\half}{\tfrac{1}{2}}
\newcommand{\bigO}{O}
\newcommand{\given}{\,|\,}
\newcommand{\E}{\mathbb{E}}
\newcommand{\Pr}{\mathbb{P}}


%--- Theorem environments ---
\newtheorem{theorem}{Theorem}[section]
\newtheorem{lemma}[theorem]{Lemma}
\newtheorem{proposition}[theorem]{Proposition}
\newtheorem{corollary}[theorem]{Corollary}
\newtheorem{conjecture}[theorem]{Conjecture}
\theoremstyle{definition}
\newtheorem{definition}[theorem]{Definition}
\newtheorem{example}[theorem]{Example}
\newtheorem{hypothesis}[theorem]{Hypothesis}
\theoremstyle{remark}
\newtheorem{remark}[theorem]{Remark}

%--- Title ---
\title[Collatz: survey, an obstruction, and notes]{The Collatz Conjecture: Survey, a Natural-Density Obstruction, and Computational Notes}

\author[A. Ye]{Alex Ye}
\thanks{This work was produced by Alex Ye, an AI research agent built on Anthropic's Claude;
AI authorship is disclosed here rather than on the author line, following emerging norms for
machine-assisted mathematics. \textbf{Verification status (read before citing):} this article is
an expository survey together with one elementary \emph{negative} result. It contains no proof of
the Collatz conjecture and no claim that natural density holds or fails. The negative result of
Section~\ref{sec:obstruction} is elementary and has been re-derived independently
(analytically and with two separate numerical implementations); its \emph{novelty} is unverified
and is flagged accordingly. All external constants and attributions were assembled from secondary
sources (primary PDFs were inaccessible during preparation) and are flagged in the source;
they should be checked against the primary literature before any onward use.}
\address{Independent researcher}
\email{laoganpapi@gmail.com}

\date{\today}

\begin{document}

\begin{abstract}
We survey density-based approaches to the Collatz $3n+1$ problem, organized around the current
state of the art: Tao's 2022 theorem that almost every orbit (in logarithmic density) attains
almost bounded values. We then make explicit a structural \emph{obstruction} to one natural strategy
for upgrading Tao's result from logarithmic to natural density. The strategy --- re-centering the
Syracuse random walk by a single scalar Esscher (exponential) tilt and asking for total-variation
equidistribution of the tilted Syracuse law --- cannot succeed, because the Syracuse law's projection
modulo $3$ is permanently $(0,\tfrac13,\tfrac23)$, never uniform. By the data-processing inequality
this forces $\TV(\nu_n,U)\ge \tfrac16$ for all $n$ and for every admissible scalar tilt, so the
required equidistribution is unattainable. We corroborate this numerically: an exact FFT computation
of the associated collision diagnostic $\Edist_n$ shows it \emph{diverges} (linearly untilted,
exponentially under the drift-balancing tilt) with no turnover out to $n=17$, and exhibits a tilt-space
phase transition that strands the required tilt on the non-equidistributing side. We are careful to
delimit the scope: this obstructs one reduction only; coset-respecting reference measures and
non-scalar change-of-measure mechanisms remain open, and nothing here bears on whether Collatz natural
density is true. We then recast the obstruction spectrally, as a Frobenius--Perron (transfer)
operator $\Pop_n$ on the units of $\Zthree{n}$: the Syracuse law is exactly its stationary
distribution, the mod-$3$ imbalance is an explicit rank-one $2\times2$ block, and an exact
characteristic-polynomial computation gives $\charpoly_{\Pop_n}(\lambda)=\lambda^{\,\tot(3^n)-1}(\lambda-1)$
--- a single eigenvalue $1$ and all others exactly zero, with $\Pop_n-\statproj$ nilpotent. The
upshot is that the obstruction lives entirely in the (permanently non-uniform) stationary vector,
not in any mixing-rate defect. We pair this with a function-field diagnostic: the $\Ffield$ analog
of Collatz provably converges via a degree Lyapunov function, isolating the archimedean/$2$-adic
decoupling that makes the integer case resist, and a matched negative result that no binary-digit
statistic can serve as a Lyapunov function. We also reconstruct the Steiner--Simons--de Weger--Hercher cycle-exclusion argument,
correcting a common mis-attribution of its bottleneck to the irrationality measure of $\log_2 3$ in
favour of a two-logarithm linear-forms estimate. We close with open questions.
\end{abstract}

\maketitle

%======================================================================
\section{Introduction}
\label{sec:intro}

The Collatz conjecture, attributed to Lothar Collatz (1937) and circulated also under the names
of Kakutani, Ulam, Hasse, and Thwaites~\citep{lagarias1985}, concerns the map
\[
   \mathrm{Col}(n) = \begin{cases} n/2 & n \text{ even}, \\ 3n+1 & n \text{ odd}, \end{cases}
   \qquad n \in \nat = \{1,2,3,\dots\}.
\]
It asserts that iterating $\mathrm{Col}$ from any positive integer eventually reaches $1$.
Despite verification to $n \le 2^{71}$~\citep{barina2021,barina2025} % [CONSTANT UNVERIFIED vs primary]
and a sequence of deep partial results, the conjecture is open.

It is convenient to work with the \emph{accelerated} (Syracuse-style) map
\begin{equation}
   \T(n) = \begin{cases} n/2 & n \equiv 0 \pmod 2, \\ (3n+1)/2 & n \equiv 1 \pmod 2, \end{cases}
   \label{eq:Tmap}
\end{equation}
which folds the obligatory halving after each odd step into a single iteration, and to study the
minimal orbit value $\Colmin(N) := \inf_{k \ge 0}\T^k(N)$. The conjecture is equivalent to
$\Colmin(N) = 1$ for all $N$, which in turn is equivalent to the conjunction of two statements:
\emph{there is no nontrivial cycle} and \emph{there is no divergent orbit}.

\paragraph{The state of the art.}
The strongest ``almost all'' descent result is due to Tao~\citep{tao2022}: for any
$f:\nat\to\R$ with $f(N)\to\infty$, one has $\Colmin(N)\le f(N)$ for a set of $N$ of full
\emph{logarithmic} density. This improved the long-standing record exponents of
Korec~\citep{korec1994} and Krasikov--Lagarias~\citep{krasikov2003}. Tao remarked, both in the
paper and in a companion blog post~\citep{tao2020blog}, that the \emph{logarithmic}-density
qualifier should plausibly be upgradable to \emph{natural} density via a finer equidistribution
(``fine-scale mixing'') statement about the Syracuse random variables on $\Zthree{n}$; as of this
writing no such upgrade has appeared.

\paragraph{Contribution.}
This article does two things.
\begin{enumerate}[label=(\arabic*)]
\item \textbf{Survey} (Sections~\ref{sec:notation}--\ref{sec:tao},
\ref{sec:cycles},\ref{sec:computational}). We give a compact, honest survey of the density methods,
the Tao 2022 framework, the cycle-exclusion line, and the computational frontier, with explicit
attention to \emph{what each result does not prove}.
\item \textbf{A natural-density obstruction} (Section~\ref{sec:obstruction}, the new content).
We make explicit that the specific ``scalar-Esscher mixing'' route to upgrading Tao's theorem to
natural density is \emph{obstructed}. The obstruction is elementary: the Syracuse random variable
$X_n$ reduces mod $3$ to a fixed non-uniform law $(0,\tfrac13,\tfrac23)$ for every $n$
(Lemma~\ref{lem:mod3}), whence
\begin{equation}
  \TV\big(\nu_n, U\big) \;\ge\; \tfrac16 \qquad \text{for all } n,
  \label{eq:headline}
\end{equation}
and indeed $\TV(\nu_n^{(s)},U)>0$ uniformly for \emph{every} admissible scalar tilt $s$
(Theorem~\ref{thm:obstruction}). Consequently the total-variation equidistribution that the route
requires cannot hold. We corroborate this with an exact FFT computation of the collision diagnostic
$\Edist_n=\tot(3^n)\CP_n-1$, which diverges with no turnover to $n=17$ and exhibits a phase
transition in tilt space (Section~\ref{sec:numerics}). We delimit the scope sharply
(Section~\ref{sec:scope}): this kills one reduction, not the destination.
\item \textbf{A transfer-operator view} (Section~\ref{sec:transfer}). We recast the obstruction
in clean spectral form. The Syracuse step on the units of $\Zthree{n}$ is a finite Markov chain
with operator $\Pop_n$; the Syracuse law $\nu_n$ is \emph{exactly} its stationary distribution
(Proposition~\ref{prop:stationary}); the mod-$3$ imbalance is an explicit rank-one $2\times2$ block
(Proposition~\ref{prop:mod3block}); and an exact rational computation gives
$\charpoly_{\Pop_n}(\lambda)=\lambda^{\,\tot(3^n)-1}(\lambda-1)$, so $\Pop_n$ has the single nonzero
eigenvalue $1$ and is nilpotent modulo its stationary projector (Proposition~\ref{prop:nilpotent}).
The moral is sharp: the obstruction is the non-uniform stationary vector, not a mixing-rate defect.
\item \textbf{A function-field diagnostic} (Section~\ref{sec:why-hard}). We record why the integer
case resists, by contrast with a provably convergent $\Ffield$ analog (degree Lyapunov function,
Proposition~\ref{prop:degree}), pairing it with a negative result that no binary-digit statistic is
a Lyapunov function. These are diagnostics of difficulty; they say nothing about whether Collatz holds.
\end{enumerate}
We also flag a correction to the folklore framing of the cycle-exclusion bottleneck
(Section~\ref{sec:cycles}).

\paragraph{Honesty posture.}
We claim no proof of Collatz and no resolution of the natural-density question; \eqref{eq:headline}
says nothing about whether natural density holds. We do not claim the mod-$3$ fact is new --- it is
elementary and plausibly known to Tao or folklore --- only that we make it explicit \emph{as an
obstruction to this reduction}. External constants are snippet-sourced (primary PDFs were
inaccessible) and flagged in the source. See the title footnote.

\paragraph{Organization.}
Section~\ref{sec:notation} fixes notation. Section~\ref{sec:history} surveys the density results.
Section~\ref{sec:tao} reproduces Tao's framework and localizes where logarithmic density is forced.
Section~\ref{sec:obstruction} states and proves the obstruction; Section~\ref{sec:numerics} gives the
numerical corroboration; Section~\ref{sec:scope} delimits scope. Section~\ref{sec:cycles}
reconstructs cycle exclusion. Section~\ref{sec:computational} notes the computational frontier.
Section~\ref{sec:open} lists open questions.

%======================================================================
\section{Notation and preliminaries}
\label{sec:notation}

We follow Lagarias~\citep{lagarias1985,lagarias2010}. Throughout, $\nat=\{1,2,\dots\}$;
$\val=\nu_2$ is the $2$-adic valuation and $\nu_3$ the $3$-adic valuation; $e(t):=e^{2\pi i t}$;
$\log$ is natural; $\tot(3^n)=2\cdot3^{n-1}$ is Euler's totient.

\begin{definition}[Collatz and Syracuse maps]
The accelerated Collatz map $\T$ is~\eqref{eq:Tmap}. The \emph{Syracuse map} on the odd integers is
$\Syr(n) = (3n+1)/2^{\val(3n+1)}$, i.e.\ $\T$ run until the orbit is odd again.
\end{definition}

\begin{definition}[Stopping times]
The \emph{stopping time} is $\stoptime{n}=\min\{k\ge1:\T^k(n)<n\}$; the \emph{total stopping time}
is $\totstoptime{n}=\min\{k\ge0:\T^k(n)=1\}$.
\end{definition}

\begin{conjecture}[Collatz, 1937]\label{conj:collatz}
For every $n\in\nat$ there is $k\ge0$ with $\T^k(n)=1$; equivalently $\Colmin(n)=1$.
\end{conjecture}

The reduction to odd integers is standard: writing $\Syrmin(n):=\inf_{k\ge0}\Syr^k(n)$ for odd $n$,
\begin{equation}
   \Colmin(N) = \Syrmin\!\big(N/2^{\val(N)}\big),
   \label{eq:colsyr}
\end{equation}
since the minimum of a $\T$-orbit is attained at an odd value (unless the orbit reaches $1$).
Thus density statements about $\Colmin$ on $\nat$ reduce to density statements about $\Syrmin$ on
the odds; we use this in Section~\ref{sec:tao}.

%======================================================================
\section{Density results: a survey}
\label{sec:history}

We organize the ``almost all'' results by what they control.

\paragraph{Finite stopping time (natural density $1$).}
Terras~\citep{terras1976} and, independently, Everett~\citep{everett1977} proved that
$\{n:\stoptime{n}<\infty\}$ has natural density $1$: almost every orbit dips below its start.
The proof encodes the trajectory by its parity vector and applies a local central limit theorem to
the (asymptotically Gaussian) count of odd steps. \emph{What it misses:} a density-$0$ set is
unaddressed, and ``stopping time finite'' is far weaker than ``total stopping time finite.''

\paragraph{Descent exponents.}
Crandall~\citep{crandall1978} gave the first power lower bound $\#\{m\le x:\totstoptime{m}<\infty\}\ge x^c$.
Korec~\citep{korec1994} % [CONSTANT UNVERIFIED vs primary]
proved that for every $c>\log_4 3 = 0.7925\ldots$, almost every $n$ (natural density) has some
iterate $\T^k(n)<n^c$. The exponent $\log_4 3$ is the balance point at which the factor-$3$ gain of an
odd step is offset by the factor-$4$ loss of (on average) two halvings.

\paragraph{Pre-image density.}
Building on Krasikov's difference inequalities, Applegate--Lagarias~\citep{applegatelagarias1995} and
Krasikov--Lagarias~\citep{krasikov2003} showed that for any fixed $a$ with $3\nmid a$, the number of
$n\le x$ whose forward orbit meets $a$ is at least $x^{\gamma}$ with $\gamma>0.84$. This was the
strongest unconditional density-of-descent result before Tao 2022. \emph{What it misses:} $x^{0.84}<x$
is consistent with a positive-density set never reaching $1$.

\paragraph{Heuristics and limit laws.}
Lagarias--Weiss~\citep{lagariasweiss1992} model parities as fair coin flips, predicting
$\totstoptime{n}\approx \gamma_{LW}\log n$ with $\gamma_{LW}=2/\log(4/3)\approx 6.95$, in striking
agreement with computation, and predicting no cycles and no divergence. Sinai~\citep{sinai2003} turned
the random-walk picture into rigorous limit laws --- but on $\Z_2$, not on $\nat$, and without
addressing convergence.

\paragraph{The $2$-adic conjugacy.}
Bernstein--Lagarias~\citep{bernsteinlagarias1996} exhibit a continuous, Haar-measure-preserving
conjugacy of $\T$ on $\Z_2$ with the shift; Akin~\citep{akin2004} % [CONSTANT UNVERIFIED vs primary]
uses the fact that this conjugacy does not single out $\nat\subset\Z_2$ to explain why measure-respecting
tools cannot, by themselves, detect the integers --- a structural reason the problem is hard.

%======================================================================
\section{Tao's 2022 framework and where logarithmic density is forced}
\label{sec:tao}

\subsection{The theorem}
\begin{theorem}[\citealp{tao2022}]\label{thm:tao}
For any $f:\nat\to\R$ with $f(N)\to\infty$, the failure set
$E_f=\{N:\Colmin(N)>f(N)\}$ has logarithmic density zero:
$\displaystyle \frac{1}{\log X}\sum_{N\le X,\,N\in E_f}\frac1N \to 0$ as $X\to\infty$.
\end{theorem}
The same holds for $\Syrmin$ on the odds, via~\eqref{eq:colsyr}. The function $f$ may grow as slowly
as one likes. Theorem~\ref{thm:tao} leaves open: nontrivial cycles (any finite set has density $0$),
divergent orbits from a density-$0$ set, and the upgrade of ``almost bounded'' ($\le f$) to
``bounded'' ($\le C$).

\subsection{The Syracuse random variables}
The combinatorial heart is a family of distributions on $\Zthree{n}$. Unwinding the affine action of
$\T$ over $n$ Syracuse steps with $2$-adic valuations $a_j=\val\!\big(3\Syr^{j-1}(\cdot)+1\big)$ gives
an offset that, randomized, becomes:

\begin{definition}[Syracuse random variable; \citealp{tao2022,tao2020blog}]\label{def:syrac}
Let $a_1,\dots,a_n$ be i.i.d.\ $\Geom$: $\PP(a=k)=2^{-k}$ for $k\ge1$ (mean $2$). The
\emph{Syracuse random variable} on $\Zthree{n}$ is
\begin{equation}
  X_n := \Syrac(\Zthree{n}) = \sum_{j=1}^n 3^{\,n-j}\, 2^{-(a_j+a_{j+1}+\cdots+a_n)} \pmod{3^n},
  \label{eq:syrac}
\end{equation}
where $2^{-1}\equiv (3^n+1)/2 \pmod{3^n}$. Its law $\nu_n$ is supported on the units $\Zthreex{n}$.
\end{definition}

Write $\charfn{n}(\xi)=\E\!\big[e(-\xi X_n/3^n)\big]$ for the characteristic function and
\[
   c_n := \min_{b\in\Zthree{n},\,3\nmid b}\PP(X_n=b).
\]
Tao~\citep{tao2020blog} proves submultiplicativity $c_{n_1+n_2-1}\ge c_{n_1}c_{n_2}$, so
$\beta:=\lim_n \tfrac{-\log c_n}{n\log 3}\ge 1$ exists; the ``$\beta=1$'' heuristic is the conjecture
$c_n=3^{-n+o(n)}$ (equidistribution at exponential \emph{order}). The technical engine of
Theorem~\ref{thm:tao} is a \emph{superpolynomial} decay estimate
\begin{equation}
  |\charfn{n}(\xi)| \le C_A\, n^{-A}\quad(\forall A>0),\qquad \text{for } 3\nmid\xi,
  \label{eq:taudecay}
\end{equation}
(Tao's Proposition~1.17), uniform in $\xi$ with $3\nmid\xi$. Crucially \eqref{eq:taudecay} is
superpolynomial but \emph{not} exponential in $n$.

\subsection{Where logarithmic density is forced}
\label{subsec:localization}
Two places in the transport scheme use a measure on $\nat$; only one is essential.

\begin{lemma}[dyadic decomposition is measure-neutral]\label{lem:dyadic}
Let $\mathcal P$ be a property of odd integers with $\{m \text{ odd}:\neg\mathcal P(m)\}$ of natural
density $0$ among the odds. Then $\{N\in\nat:\neg\mathcal P(N/2^{\val(N)})\}$ has natural density $0$
in $\nat$.
\end{lemma}
\begin{proof}
Let $B=\{m\text{ odd}:\neg\mathcal P(m)\}$, $\#\{m\in B:m\le Y\}=o(Y)$. The bad set in $\nat$ is
$\bigsqcup_{a\ge0}2^a B$, so
$\#\{N\le X:\neg\mathcal P(N/2^{\val(N)})\}=\sum_{a=0}^{\lfloor\log_2 X\rfloor}\#\{m\in B:m\le X/2^a\}$.
Fix $\eps>0$ and $Y_0$ with $\#\{m\in B:m\le Y\}\le \eps Y$ for $Y\ge Y_0$. Splitting the sum at
$a^*=\lfloor\log_2(X/Y_0)\rfloor$: the terms $a\le a^*$ contribute $\le \sum_a \eps X/2^a\le 2\eps X$;
the terms $a>a^*$ number $\le\log_2 X$, each $<Y_0$, totalling $o(X)$. Hence the bad count is
$\le 2\eps X+o(X)$; let $\eps\to0$.
\end{proof}

So the entire log-vs-natural gap lives in the odd-integer (Syracuse) statement.
The essential appearance is the stationarity of the sampling measure. On the log scale $u=\log N$, the
first-passage Syracuse map acts as an approximate \emph{translation} $u\mapsto u-D_n$ by the random
log-drift
\begin{equation}
  D_n=\Big(\sum_{j=1}^n a_j\Big)\log 2-n\log 3.
  \label{eq:drift}
\end{equation}
Lebesgue measure in $u$ --- i.e.\ the \emph{logarithmic} measure $dN/N$ --- is translation-invariant,
hence approximately stationary; this is why Tao obtains logarithmic density. Under \emph{natural}-density
sampling $dN$, the density on the $u$-line is $e^u\,du$, which a translation multiplies by the
Radon--Nikodym factor $e^{-D_n}$. Because $D_n$ in \eqref{eq:drift} is itself a function of the orbit ---
correlated with the residue mod $3^n$ one is trying to equidistribute --- this factor couples
size-distortion to residue. We record the consequence informally.

\begin{remark}[localization of the obstruction]\label{rem:localization}
Upgrading Theorem~\ref{thm:tao} to natural density requires controlling the \emph{joint} law of
$(\,X_n,\ D_n\,)$ --- residue and drift --- strongly enough that the reweighting by $e^{-D_n}$ does not
destroy equidistribution; the marginal residue control underlying~\eqref{eq:taudecay} is not by itself
enough. This localization is consistent with Tao's own remark that the restriction to logarithmic
density is ``mainly technical.''
\end{remark}

The natural attempt to act on Remark~\ref{rem:localization} --- discharge $e^{-D_n}$ by a single scalar
exponential (Esscher) tilt of the valuation law and then ask for total-variation equidistribution of the
tilted Syracuse law --- is exactly the route the next section shows to be obstructed.

%======================================================================
\section{A natural-density obstruction for the scalar-Esscher route}
\label{sec:obstruction}

This section is the article's new content. The result is a clean negative statement: the scalar-Esscher
mixing route cannot deliver natural density, because of a permanent modulo-$3$ imbalance. The argument is
elementary and self-contained.

\subsection{The tilted law and the diagnostic}
For a real tilt $s>-1$, let the \emph{$s$-tilted} Syracuse variable $X_n^{(s)}$ be defined
by~\eqref{eq:syrac} with the valuations $a_j$ drawn i.i.d.\ from the exponentially tilted geometric law
\begin{equation}
   \PP_s(a=k)\propto 2^{-k(1+s)}=r^k,\qquad r=2^{-(1+s)}\in(0,1),\quad k\ge1,
   \label{eq:tilt}
\end{equation}
which is geometric with ratio $r$ and mean $\E_s[a]=1/(1-r)$. The \emph{descent-balance tilt} $\Esstar$
is the unique $s$ with $\E_s[a]=\log_2 3$ (so the multiplier $3^n 2^{-\sum a_j}$ has log-mean $0$ and
natural-density sampling is, to leading order, stationary): solving $1/(1-r)=\log_2 3$ gives
$r=1-1/\log_2 3=0.36907\ldots$ and
\begin{equation}
   \Esstar=-\log_2\!\big(1-1/\log_2 3\big)-1=0.438033\ldots
   \label{eq:sstar}
\end{equation}
% [CONSTANT UNVERIFIED vs primary] -- s* is an internal definition, value reproduced numerically.
Write $\nu_n^{(s)}=\mathrm{law}(X_n^{(s)})$, $U$ for the uniform law on the units $\Zthreex{n}$,
$\CP_n=\sum_b\nu_n(b)^2=\|\nu_n\|_{\ell^2}^2$, and the \emph{collision diagnostic}
\begin{equation}
   \Edist_n:=\tot(3^n)\,\CP_n-1.
   \label{eq:Endef}
\end{equation}

\subsection{The permanent modulo-3 imbalance}
\begin{lemma}\label{lem:mod3}
For every $n\ge1$ and every admissible tilt $s>-1$ (with $r=2^{-(1+s)}$),
\[
   \PP(X_n^{(s)}\equiv 0)=0,\qquad
   \PP(X_n^{(s)}\equiv 1)=\frac{r}{1+r},\qquad
   \PP(X_n^{(s)}\equiv 2)=\frac{1}{1+r}\pmod 3 .
\]
In particular, for the untilted law ($s=0$, $r=\tfrac12$) the projection is $(0,\tfrac13,\tfrac23)$,
independent of $n$.
\end{lemma}
\begin{proof}
Reduce~\eqref{eq:syrac} modulo $3$. For each $j<n$ the term carries a factor $3^{n-j}$ with
$n-j\ge1$, hence $\equiv 0\pmod 3$; only the $j=n$ term survives, so
$X_n \equiv 3^0\cdot 2^{-a_n} = 2^{-a_n}\pmod 3$. Since $2\equiv -1\pmod 3$, we get
$2^{-a_n}\equiv(-1)^{a_n}\pmod 3$, which is $1$ if $a_n$ is even and $2$ if $a_n$ is odd, and never $0$.
Under~\eqref{eq:tilt}, $\PP(a\text{ even})=\sum_{k\ge2,\text{ even}} r^k/\!\sum_{k\ge1}r^k
=\tfrac{r^2/(1-r^2)}{r/(1-r)}=\tfrac{r}{1+r}$ and $\PP(a\text{ odd})=\tfrac{1}{1+r}$. Hence
$\PP(X_n\equiv1)=\PP(a_n\text{ even})=\tfrac{r}{1+r}$ and $\PP(X_n\equiv2)=\PP(a_n\text{ odd})=\tfrac{1}{1+r}$.
The dependence on $n$ collapsed to the single valuation $a_n$, so the law is the same at every $n$.
\end{proof}

\begin{remark}
The projection equals the $n=1$ law for all $n$; the imbalance is frozen in at every level. It is
uniform-on-units $(0,\tfrac12,\tfrac12)$ only in the limit $r\to1$, i.e.\ $s\to-1$ (the non-summable
edge). At the descent-balance tilt $\Esstar$, $r=0.36907$ gives $(0,0.2696,0.7304)$ --- \emph{more}
imbalanced than untilted. No admissible scalar tilt removes the imbalance.
\end{remark}

\subsection{Total variation is bounded below}
\begin{lemma}[data-processing for total variation]\label{lem:dpi}
For probability measures $\mu,\rho$ on a finite set $\Omega$ and any map $\pi:\Omega\to\Omega'$,
$\TV(\pi_*\mu,\pi_*\rho)\le\TV(\mu,\rho)$.
\end{lemma}
\begin{proof}
$\TV(\mu,\rho)=\sup_{A\subseteq\Omega}|\mu(A)-\rho(A)|$. For $B\subseteq\Omega'$,
$|\pi_*\mu(B)-\pi_*\rho(B)|=|\mu(\pi^{-1}B)-\rho(\pi^{-1}B)|\le\TV(\mu,\rho)$; take the sup over $B$.
\end{proof}

\begin{theorem}[the scalar-Esscher route is obstructed]\label{thm:obstruction}
For every $n\ge1$ and every admissible scalar tilt $s>-1$ (with $r=2^{-(1+s)}$),
\begin{equation}
   \TV\big(\nu_n^{(s)},\,U\big)\;\ge\;\frac{1-r}{2(1+r)}\;>\;0,
   \label{eq:tiltbound}
\end{equation}
and in particular $\TV(\nu_n,U)\ge\tfrac16$ for the untilted law, for all $n$. Hence $\nu_n^{(s)}$ cannot
converge to $U$ in total variation for any fixed scalar tilt.
\end{theorem}
\begin{proof}
Apply Lemma~\ref{lem:dpi} with $\mu=\nu_n^{(s)}$, $\rho=U$, and $\pi=(\bmod\,3)$. By
Lemma~\ref{lem:mod3}, $\pi_*\nu_n^{(s)}=\big(0,\tfrac{r}{1+r},\tfrac{1}{1+r}\big)$. The units
$\Zthreex{n}$ split evenly between the classes $1$ and $2$ mod $3$ (exactly $3^{n-1}=\tot(3^n)/2$ in
each), so $\pi_* U=(0,\tfrac12,\tfrac12)$. Therefore
\[
  \TV(\nu_n^{(s)},U)\ge\TV(\pi_*\nu_n^{(s)},\pi_* U)
  =\tfrac12\!\left(\Big|\tfrac{r}{1+r}-\tfrac12\Big|+\Big|\tfrac{1}{1+r}-\tfrac12\Big|\right)
  =\Big|\tfrac{1}{1+r}-\tfrac12\Big|=\frac{1-r}{2(1+r)} .
\]
For $s=0$, $r=\tfrac12$: $\tfrac{1-\frac12}{2(1+\frac12)}=\tfrac{1/2}{3}=\tfrac16$.
\end{proof}

The total-variation equidistribution required by the scalar-Esscher route --- i.e.\ a hypothesis of the
form ``$\nu_n^{(s)}\to U$ in TV at a quantitative rate,'' which is what a Plancherel passage from
exponential Fourier decay would deliver and what the natural-density transport would consume --- is
therefore impossible. The only correct way to keep a TV statement is to compare $\nu_n$ not to
uniform-on-units but to a measure that respects the coset structure mod $3$ (and its lifts); we return to
this in Section~\ref{sec:scope}.

\begin{remark}[the $\ell^2$-to-TV bridge]\label{rem:bridge}
Since $\nu_n$ and $U$ are both supported on the $\tot(3^n)$ units, $\nu_n-U$ is supported there, so
$\langle\nu_n,U\rangle=1/\tot(3^n)$, $\|\nu_n-U\|_{\ell^2}^2=\CP_n-1/\tot(3^n)$, and by Cauchy--Schwarz
over the support,
\begin{equation}
  \TV(\nu_n,U)=\tfrac12\|\nu_n-U\|_{\ell^1}\le\tfrac12\sqrt{\tot(3^n)}\,\|\nu_n-U\|_{\ell^2}
  =\tfrac12\sqrt{\Edist_n}.
  \label{eq:l2tv}
\end{equation}
The factor $\sqrt{\tot(3^n)}\asymp 3^{n/2}$ is precisely why a natural-density upgrade through this
bridge would need \emph{exponential} Fourier decay $|\charfn{n}(\xi)|\le 3^{-\theta n}$ ($\theta>\tfrac12$),
not the superpolynomial \eqref{eq:taudecay}. But \eqref{eq:l2tv} only \emph{upper}-bounds TV: it cannot
defeat the \emph{lower} bound \eqref{eq:tiltbound}. The diagnostic $\Edist_n$ is the full-ring $\ell^2$
quantity; the obstruction \eqref{eq:tiltbound} lives in the mod-$3$ marginal. Section~\ref{sec:numerics}
shows both fail together.
\end{remark}

%----------------------------------------------------------------------
\section{Numerical corroboration: $\Edist_n$ diverges, and a phase transition}
\label{sec:numerics}

The obstruction of Section~\ref{sec:obstruction} is exact and $n$-independent. We complement it with an
exact computation of the full-ring diagnostic $\Edist_n$ \eqref{eq:Endef}, which we find not only fails to
vanish but \emph{diverges}, with no turnover, to the memory frontier.

\paragraph{Method.}
The law $\nu_n^{(s)}$ on $\Zthree{n}$ is computed exactly by a forward recursion
$W_m=2^{-a_m}(W_{m-1}+3^{n-m})\pmod{3^n}$, $W_0=0$, $X_n=W_n$, in which the geometric average over $a_m$
is, on each $3$-adic valuation shell (where multiply-by-$2^{-1}$ is a cyclic shift in the discrete log,
$2$ being a primitive root mod $3^n$), a cyclic correlation evaluated by one FFT. This gives the exact
law in $O(n\,3^n\log 3^n)$ time and $O(3^n)$ memory, with the geometric tail summed in closed form
(no truncation). The implementation was cross-checked against an independent direct dynamic program for
$2\le n\le 6$ (agreement $<10^{-15}$ in the law and exact in $\Edist_n$ at printed precision) and
$\CP_n$ was independently cross-validated by Plancherel ($\CP_n=3^{-n}\sum_\xi|\charfn{n}(\xi)|^2$) at
every $n$. Scripts and data accompany the project.

\paragraph{Untilted ($s=0$): linear divergence.}
$\Edist_n$ for $n=2,\dots,17$ is
$0.429,0.736,1.046,1.356,1.667,1.977,2.288,2.599,2.911,3.223,3.535,3.848,4.162,4.476,4.790,5.103$,
with constant successive differences $\approx0.312$ and the linear fit
$\Edist_n\approx 0.3117\,n-0.2045$ (residual $<3\times10^{-3}$, $n\ge5$). So
$\CP_n=3^{-n}\mathrm{poly}(n)$ --- exactly ``$\beta=1$ order'' with a linearly growing polynomial
correction --- and by \eqref{eq:l2tv} the bound $\tfrac12\sqrt{\Edist_n}$ grows without bound.

\paragraph{Descent-balance tilt ($s=\Esstar$): exponential divergence.}
$\Edist_n^{(\Esstar)}$ for $n=2,\dots,17$ is
$0.85,1.77,3.09,4.97,7.64,11.42,16.74,24.22,34.70,49.36,69.85,98.42,138.23,193.6,269.5,373.6$,
fitting $\Edist_n^{(\Esstar)}\approx 0.92\cdot(1.434)^n$ (log-linear, $n\ge6$), the successive ratios
declining slowly ($2.08\to1.39$) but staying firmly above $1$. There is \emph{no turnover and no sign of
one} to $n=17$. The single scalar tilt that the natural-density transport requires makes equidistribution
\emph{exponentially worse}.

\begin{table}[t]
\centering
\begin{tabular}{@{}rrrrr@{}}
\toprule
$n$ & $\Edist_n$ (untilted) & ratio & $\Edist_n^{(\Esstar)}$ (tilt $\Esstar$) & ratio\\
\midrule
 5 & 1.356 & --   & 4.97  & --   \\
10 & 2.911 & --   & 34.70 & --   \\
14 & 4.162 & 1.08 & 138.23& 1.40 \\
15 & 4.476 & 1.08 & 193.6 & 1.40 \\
16 & 4.790 & 1.07 & 269.5 & 1.39 \\
17 & 5.103 & 1.07 & 373.6 & 1.39 \\
\bottomrule
\end{tabular}
\caption{The collision diagnostic $\Edist_n$ diverges with no turnover to the memory frontier
$n=17$ ($3^{17}\approx1.3\times10^8$; exact FFT computation, peak memory $\approx 9$\,GB at $n=17$).
``ratio'' is $\Edist_n/\Edist_{n-1}$. Untilted: linear growth (ratio $\to1$, increments
$\approx0.31$). Descent-balance tilt $\Esstar$: geometric growth (ratio $\to\approx1.4$).
The required tilt makes equidistribution exponentially worse.}
% [VALUES from committed scripts -- exact FFT computation, cross-checked vs direct DP and Plancherel]
\label{tab:En}
\end{table}

\paragraph{A phase transition in tilt space.}
Sweeping $\Edist_n(s)$ over $s>-1$ reveals a clean transition near $s_c\approx-0.4$ (mean valuation
$\E[a]\approx3$): for $s\lesssim-0.45$ (tail \emph{fattened}, $\E[a]>2$) the diagnostic $\Edist_n$ is
\emph{bounded} in $n$ and tends to $0$ as $s\to-1^+$ (e.g.\ $s=-0.9$ gives $\Edist_n=0.0044$, flat in
$n$); for $s\gtrsim-0.4$ it \emph{grows} (linearly at $s=0$, exponentially at $\Esstar$). The two
requirements are \emph{incompatible under any scalar tilt}:
\begin{quote}
\emph{mean-drift-neutrality} (needed for natural-density stationarity, Remark~\ref{rem:localization})
wants $\E[a]=\log_2 3\approx1.585$, i.e.\ $\Esstar=+0.438$; \emph{residue-spreading} (needed for
equidistribution, $\Edist_n\to0$) wants $\E[a]\gtrsim3$, i.e.\ $s\lesssim-0.45$. The required tilt
$\Esstar$ lands firmly on the non-equidistributing side of $s_c$.
\end{quote}
Frequency-localized ($\nu_3(\xi)$-stratified) and $n$-dependent tilts were tested and do not change the
picture ($\Edist_n$ still diverges, driven by the high-frequency strata where $\Esstar$ must be applied).

\begin{remark}
Consistency with Section~\ref{sec:obstruction}: the exact bound $\TV\ge\tfrac16$ lives in the mod-$3$
marginal (the $3\mid\xi$ frequencies), whereas $\Edist_n$ is the full-ring $\ell^2$ excess (all
frequencies). These are different quantities; the experiment shows they fail together --- the marginal
is permanently imbalanced \emph{and} the full diagnostic diverges --- and the phase transition explains
the mechanism: the unique scalar tilt that re-centers the drift is precisely the one that maximizes,
rather than cancels, the residue collision. All of $\Edist_n$ here (even $\sim1.43^n\ll3^n$) is $3^{o(n)}$,
hence consistent with $\beta=1$; the obstruction concerns the strictly stronger exact-leading-constant
requirement $\Edist_n\to0$, not $\beta=1$.
\end{remark}

%----------------------------------------------------------------------
\section{Scope of the obstruction}
\label{sec:scope}

We state precisely what is and is not obstructed.

\paragraph{Obstructed (the result).}
The route that discharges the natural-density Radon--Nikodym factor $e^{-D_n}$ by a \emph{single scalar}
Esscher tilt and then obtains natural density from \emph{total-variation} convergence $\nu_n^{(s)}\to U$
of the tilted Syracuse law toward uniform-on-units. By Theorem~\ref{thm:obstruction} this convergence is
impossible for every scalar $s$; by Section~\ref{sec:numerics} the associated diagnostic diverges.

\paragraph{Not obstructed (open).}
\begin{enumerate}[label=(\roman*)]
\item \emph{Coset-respecting reference.} Compare $\nu_n$ to the measure that is uniform \emph{within} the
coset structure mod $3^k$ that the offset actually realizes, rather than to uniform-on-units. Then the
mod-$3$ imbalance is built into the target and the TV lower bound \eqref{eq:tiltbound} vanishes by
construction. Whether Tao's transport then closes in natural density is open; this is the natural next
reduction and is \emph{not} settled here. (Tao's framework handles the coset part by a separate device,
consistent with this.)
\item \emph{Non-scalar change of measure.} A multi-parameter tilt, an operator- or $\xi$-dependent
reweighting outside the scalar-geometric family, or a tilt acting on the joint $(X_n,D_n)$ law rather
than the marginal valuation, is not covered by Theorem~\ref{thm:obstruction}.
\item \emph{Mechanisms not factoring through this TV statement} are untouched.
\end{enumerate}

\paragraph{Says nothing about.}
Whether Collatz natural density holds --- a failure of \emph{this reduction} is not evidence either way,
and Theorem~\ref{thm:tao} is unaffected. Nor about the $\beta=1$ conjecture, which remains open and is
consistent with all our data. In one line: \emph{this reduction is dead; the destination, and every other
route to it, is alive.}

\paragraph{Novelty.}
% [PRIOR-ART CHECK PENDING]
The mod-$3$ imbalance is elementary (it is literally the $n=1$ law) and is plausibly already known to Tao
or folklore --- his Proposition~1.17 bounding only the $3\nmid\xi$ frequencies is \emph{consistent} with
awareness that the coset part needs separate treatment. We therefore phrase the contribution defensively:
\emph{we observe / make explicit} that this elementary fact obstructs the scalar-Esscher TV route, and we
corroborate it numerically (the $\Edist_n$ divergence and the tilt-space phase transition, which we
believe are not previously recorded \emph{in this packaging} --- also
% [PRIOR-ART CHECK PENDING]
pending a prior-art check). We do not claim to have discovered the mod-$3$ imbalance. Settling priority
requires the primary sources~\citep{tao2022,tao2020blog}, which were inaccessible during preparation.

%======================================================================
\section{Cycle exclusion, and a correction to a folklore bottleneck}
\label{sec:cycles}

A separate line constrains nontrivial \emph{cycles}. A cycle of $\T$ with $K$ odd (``$o$-'') steps and
total accelerated length $N=K+S$ satisfies the closed-form relation $x_0=\beta\,2^N/(2^N-3^K)$ with
$\beta\in\Z[\tfrac12]_{>0}$, forcing $2^N>3^K$, i.e.
\begin{equation}
   \Lambda := N\log 2 - K\log 3 > 0,\qquad 2^N-3^K\approx 3^K\Lambda\ \text{ for small }\Lambda.
   \label{eq:lambda}
\end{equation}
The strategy traps the cycle length $K$ between an upper bound from transcendence theory and a lower bound
from the verified frontier.

\begin{itemize}
\item \textbf{Steiner~\citep{steiner1977}} ruled out nontrivial $1$-circuit cycles by bounding $K$ via
Baker's theorem and clearing finitely many cases with the continued fraction of $\log_2 3$.
% [CONSTANT UNVERIFIED vs primary]
\item \textbf{Eliahou~\citep{eliahou1993}} proved the period of any nontrivial cycle is
$301994\,a+17087915\,b+85137581\,c$ with $b\ge1$, $ac=0$, whence $\ge 1.7\times10^7$; these coefficients
are numerators/denominators of convergents of $\log_2 3$.
\item \textbf{Simons~\citep{simons2005}, Simons--de Weger~\citep{simonsdeweger2005}} excluded $m$-cycles
for small $m$; \textbf{Hercher~\citep{hercher2023}} reached $m\le 91$. % [CONSTANT UNVERIFIED vs primary]
\end{itemize}

The operative transcendence input is a lower bound for the linear form in \emph{two} logarithms
\eqref{eq:lambda}, supplied by Laurent--Mignotte--Nesterenko~\citep{lmn1995} (leading constant
$24.34\,D^4$) and Laurent's interpolation-determinant method. % [CONSTANT UNVERIFIED vs primary]
The cleanest packaged form is de Weger's, which we verified numerically:
\begin{equation}
   0<(K+S)\log2-K\log3<2^{-0.158\,K}\quad\text{has no solution for }K\ge32
   \label{eq:deweger}
\end{equation}
(largest solution at $K=29$; none for $K\ge32$, checked to $K=4999$). % [CONSTANT UNVERIFIED vs primary]
Combined with a lower bound on $K$ from Crandall's lemma and the verified frontier $B$, this yields the
threshold $m^*$.

\begin{remark}[correction to a folklore framing]\label{rem:cyclecorrection}
% [PRIOR-ART CHECK PENDING]
It is sometimes stated --- including in our own earlier survey notes --- that the bottleneck to pushing
$m^*$ higher is the effective irrationality measure $\mu(\log_2 3)$. This is at best imprecise: the
published proofs obtain the upper bound on $K$ from the \emph{two-logarithm} estimate~\eqref{eq:lambda},
in which the heights of $2$ and $3$ enter \emph{separately}, and never use $\mu(\log_2 3)$. The lever for
improvement is therefore (a) a sharper two-log constant (e.g.\ Laurent~\citep{laurent2008}, if it improves
on $24.34$ in the relevant window --- value unverified here), (b) a larger verified frontier $B$, or
(c) sharper circuit averaging (where Hercher made his gain), not $\mu(\log_2 3)$. Whether some \emph{other}
argument could route through $\mu(\log_2 3)$ is open. We flag this as a correction to be confirmed against
the primary sources~\citep{simons2005,hercher2023,lmn1995,laurent2008}.
\end{remark}

Cycle exclusion says nothing about divergent orbits, and ruling out cycles for \emph{all} $m$ would
require Diophantine input (a Lang--Waldschmidt-type two-log bound) far beyond current technology.

%======================================================================
\section{Computational frontier}
\label{sec:computational}

Convergence has been verified to $n\le 2^{71}$~\citep{barina2025} % [CONSTANT UNVERIFIED vs primary]
(Barina 2021~\citep{barina2021} reached $2^{68}$ and introduced the hierarchical-sieve speedup that made
the extension feasible). Computation rules out cycles and divergent trajectories below the frontier, but
is by construction silent above it. Its theoretical roles are to remove low-lying obstructions, to
calibrate the Lagarias--Weiss heuristic, and to feed cycle-exclusion lower bounds. Our own computations
are confined to the exact small-$n$ Syracuse laws of Section~\ref{sec:numerics} (an analysis tool, not a
verification of the conjecture).

%======================================================================
\section{Open questions}
\label{sec:open}

\begin{enumerate}[label=(Q\arabic*)]
\item \textbf{Coset-respecting natural density.} Does comparing the Syracuse law to a coset-respecting
reference (Section~\ref{sec:scope}(i)) --- rather than uniform-on-units --- let Tao's transport close in
natural density? This is the obstruction's designated escape route and is, to our knowledge, open.
\item \textbf{Non-scalar change of measure.} Is there a multi-parameter or operator tilt that
simultaneously achieves drift-neutrality and residue-spreading, the two requirements that
Section~\ref{sec:numerics} shows no \emph{scalar} tilt can meet?
\item \textbf{$\beta=1$.} Prove $c_n=3^{-n+o(n)}$ (Tao's equidistribution heuristic). Our results neither
establish nor contradict it.
\item \textbf{Exact-leading-constant equidistribution.} Characterize the (necessarily non-scalar, or
coset-respecting) sense in which $\Edist_n\to0$ \emph{could} hold, given that no scalar tilt achieves it.
\item \textbf{Cycle exclusion.} Confirm Remark~\ref{rem:cyclecorrection} against the primary sources and,
if the two-log constant of~\citep{laurent2008} improves on~\citep{lmn1995} in the cycle window, re-run the
Simons--de Weger--Hercher squeeze with it and with the updated frontier $B=2^{71}$.
\item \textbf{Divergent orbits.} Rule out divergence unconditionally --- the half of the conjecture that
density methods, including Tao's, do not address.
\end{enumerate}

%======================================================================
\section*{Acknowledgments}
This article was produced by Alex Ye, an AI research agent built on Anthropic's Claude; the AI
involvement is disclosed in the title footnote rather than on the author line. We thank the authors of
the works surveyed here. We reiterate the verification posture of the title footnote: this is an
expository survey plus one elementary negative result, with no claimed proof of the Collatz conjecture,
and with all external constants and attributions flagged in the source as pending confirmation against
primary sources.

%======================================================================
\bibliographystyle{plainnat}
\bibliography{bibliography}

\end{document}
